\subsection{微分算子}

设 \(n \in N\) ， U 是 \(R^{n}\) 的开子集。我们给 \(R^{n}\) 与 U 赋以勒贝格测度，记作 \(\mu\) 。

设 \(k \in N\) 与 \(h \in N\) 满足 \(h \geqslant k\) 。设 D 是 U 上阶 \(\leqslant k\) 的纯量微分算子，其系数 \((n_{\alpha})_{|\alpha| \leqslant k}\) 在 U 上属于类 \(C^{h}\) 。假设对每个满足 \(\alpha\) 的 \(|\alpha| \leqslant k\) 以及每个满足 \(\beta\) 的 \(0 \leqslant \beta \leqslant \alpha\) ，函数 \(\partial^{\beta} n_{\alpha}\) 在 U 上有界。

设 \({}^{t}D\) 是 D 的转置纯量微分算子（VAR, R2, 14.3.2）；它的阶 \(\leqslant k\) ，属于类 \(C^{h-k}\) ；对 \(\varphi\in\mathcal{D}(U)\) ，我们有

\[
{ } ^ { t } \mathrm{D} ( \varphi ) = \sum _ { | \alpha | \leqslant k } ( - 1 ) ^ { | \alpha | } \partial ^ { \alpha } ( \overline { { n } } _ { \alpha } \varphi )
\]

（loc. cit.）；特别地， \({}^{t}D\) 的系数在 U 上有界。

我们用 D\_ 表示 \(L^{2}(U)\) 上以 \(\mathcal{D}(U)\) 为定义域、由下式定义的部分算子：

\[
\varphi \mapsto \mathrm{D} (\varphi) = \sum_ {| \alpha | \leqslant k} n _ {\alpha} \partial^ {\alpha} \varphi .\tag{5}
\]

设 \(H_{D}\) 是这样的 \(f \in L^{2}(U)\) 所成的空间：下列广义函数

\[
\mathrm{D} (f) = \sum_ {| \alpha | \leqslant k} n _ {\alpha} \partial^ {\alpha} f
\]

属于 \(L^{2}(U)\) ；我们用 \(D_{+}\) 表示以 \(H_{D}\) 为定义域、由 \(f \mapsto \mathrm{D}(f)\) 所定义的部分算子。

由于我们有 \(\partial^{\alpha}f\in \mathrm{L}^{2}(\mathrm{U})\) （若 \(f\in \mathrm{H}^k (\mathrm{U})\) 且 \(|\alpha |\leqslant k\) ），索伯列夫空间 \(\mathrm{H}^k (\mathrm{U})\) 包含在 \(\mathrm{H}_{\mathrm{D}}\) 中；一般而言，这两个空间是不同的。

我们有 \(D_{-} \subset D_{+}\) ，它们都是 D 在 \(L^{2}(U)\) 上的相伴微分算子（IV, p. 235 的定义 5）。

命题 13. --- 设 u 是 L²(U) 上的部分算子。若 D₋ ⊂ u ⊂ D₊ ，则 u 可闭，且 (ᵗD)₋ ⊂ u* ⊂ (ᵗD)₊ 。

设 \(\varphi\) 与 \(\psi\) 属于 \(\mathcal{D}(\mathrm{U})\) 。这时我们有

\[
\begin{array}{c} \langle \varphi | \mathrm{D} (\psi) \rangle = \sum_ {| \alpha | \leqslant k} \int_ {\mathrm{U}} n _ {\alpha} \overline {{\varphi}} \partial^ {\alpha} \psi d \mu \\ = \sum_ {| \alpha | \leqslant k} (- 1) ^ {| \alpha |} \langle \partial^ {\alpha} (\overline {{n}} _ {\alpha} \varphi) | \psi \rangle = \langle {} ^ {t} \mathrm{D} (\varphi) | \psi \rangle \end{array}
\]

（参见 VAR, R2, 14.3.8）。由于 \(\mathcal{D}(\mathrm{U})\) 在 \(\mathrm{L}^2 (\mathrm{U})\) 中稠密，这蕴含 \(\varphi \in \mathrm{dom}(u^{*})\) 且 \(u^{*}(\varphi) = {}^{t}\mathrm{D}(\varphi)\) 。因此我们有 \(({}^{t}\mathrm{D})_{-}\subset u^{*}\) 。特别地， \(u^{*}\) 稠定而 \(u\) 可闭（IV, p. 237 的命题 8）。

设 \(f \in \mathrm{dom}(u^{*})\) 与 \(\varphi \in \mathcal{D}(\mathrm{U})\) 。由于 \(\mathcal{D}(\mathrm{U}) \subset \mathrm{dom}(u)\) ，与 \(u^{*}(f)\) 相伴的广义函数满足

\[
\langle u ^ {*} (f), \varphi \rangle = \langle \overline {{\varphi}} | u ^ {*} (f) \rangle = \langle u (\overline {{\varphi}}) | f \rangle .
\]

由于 \(D_{-} \subset u\) ，可由公式 (5) 计算 \(u(\overline{\varphi})\) ，由此

\[
\begin{array}{l} \langle u ^ {*} (f), \varphi \rangle = \sum_ {\alpha} \langle n _ {\alpha} \partial^ {\alpha} \overline {{\varphi}} | f \rangle = \sum_ {\alpha} \langle \partial^ {\alpha} \overline {{\varphi}} | \overline {{n}} _ {\alpha} f \rangle \\ \qquad = \sum_ {\alpha} \langle \overline {{n}} _ {\alpha} f, \partial^ {\alpha} \varphi \rangle = \sum_ {\alpha} (- 1) ^ {| \alpha |} \langle \partial^ {\alpha} (\overline {{n}} _ {\alpha} f), \varphi \rangle = \langle {} ^ {t} D (f), \varphi \rangle . \end{array}
\]

于是广义函数 \(u^{*}(f)\) 与 \({}^{t}\mathrm{D}(f)\) 相等；因而广义函数 \(f\) 属于 \(H_{tD}\) ，且 \(u^{*}(f)={}^{t}\mathrm{D}(f)\) ，由此 \(u^{*}\subset({}^{t}\mathrm{D})_{+}\) 。

注. --- 这一命题表明，部分算子 \(u\) （它满足 \(\mathrm{D}_{-} \subset u \subset \mathrm{D}_{+}\) ）的伴随总可以在广义函数的意义下计算：元素 \(f\) （它们属于 \(u^{*}\) 的定义域）是 \(\mathrm{L}^{2}(\mathrm{U})\) 中这样的元素：广义函数 \(^{t}\mathrm{D}(f)\) 属于 \(\mathrm{L}^{2}(\mathrm{U})\) ，且我们有 \(u^{*}(f) = ^{t}\mathrm{D}(f)\) 。

若作为纯量微分算子有 \({}^{t}D = D\) ，就称 D 是形式对称的。在这种情形，部分算子 D\_ 是对称的。

考虑由下式定义的二阶纯量微分算子的特殊情形：

\[
\Delta = - \sum_ {i = 1} ^ {n} \partial_ {i} ^ {2}.
\]

拉普拉斯算子（ U 上的），指的是 \(L^{2}(U)\) 上自伴的部分算子 u ，满足 \(\Delta_{-} \subset u\) （参见 VAR, R2, 14.4.3, p. 83）。我们稍后将看到（IV, p. 261，例）， \(L^{2}(U)\) 上总存在至少一个拉普拉斯算子；也可能存在不止一个（IV, p. 358 的习题 17）。
% semantic-fragments: legacy-input-seam
\relax{}
